%% Согласно ГОСТ Р 7.0.11-2011:
%% 5.3.3 В заключении диссертации излагают итоги выполненного исследования, рекомендации, перспективы дальнейшей разработки темы.
%% 9.2.3 В заключении автореферата диссертации излагают итоги данного исследования, рекомендации и перспективы дальнейшей разработки темы.
\begin{enumerate}
  \item На основе анализа данных интернет-трафика от Канадского Института Кибербезопасности был проведен углубленный анализ ранее предложенных решений за пределами оригинальных условий. Получилось расширить область применения метода AGMV и получить прирост метрики accuracy не менее чем на 0,2, а для метода CAPoNeF выявить границы его применимости на основе экспериментального анализа. Полученные выводы могут быть использованы для более осознанного выбора стратегий генерации признаков.
  
  \item На основе данных RIPE Atlas разработаны методы предобработки и калибровки показателей времени отклика, учитывающие особенности устройств и временные флуктуации. Это позволило повысить точность оценки производительности узлов путем устранения искажений у 63\% используемых устройств.
  
  \item На основе анализа данных о RTT, собранных с измерительных узлов, расположенных на территории Российской Федерации, был разработан набор методов выявления аномальных устройств, демонстрирующих отклонения от типового сетевого поведения. Применение разработанного метода позволило выделить 32\% аномальных узлов, фильтрация которых привела к снижению ошибки (RMSE) оценки времени приема передачи до 50\%.
  
  \item На основе обработанных данных был разработан подход к оценке быстродействия сетевых узлов. Предложенный метод требует на 3–4 порядка меньше операций с плавающей запятой (FLOPs), что делает его пригодным для масштабирования и применения в условиях ограниченных вычислительных ресурсов.
\end{enumerate}
